\begin{center}
\LARGE
Semantic Pattern Matching In Mathematica
\end{center}

\begin{center}
\Large
Jason Harris

\end{center}

\noindent
Date:  July 19th (Friday)\\
Time:  09:20-09:45\\

\begin{center}
\large
Abstract
\end{center}

This talk presents a novel extension to the Mathematica language to allow semantic pattern
matching. We present examples of semantic pattern matching that demonstrate the full range of
abilities of this semantic extension. Semantic patterns function correctly with single patterns, free
symbols, structural symbols, constrained patterns in particular PatternTests, Conditions and
Optionals, nonlinear semantic patterns, named pattern objects, assignments and transformation
rules. In short semantic patterns uniformly extended the normal pattern matching mechanisms in
Mathematica. 

Semantica compiles semantic patterns into corresponding syntactic patterns, which then function
at the normal speed of Mathematica. Semantic patterns integrate rewriting, algebraic unification,
and equational solving in an intuitive and seamless way. 

